\documentclass[12pt]{article}
\usepackage[utf8]{inputenc}
\usepackage[T1]{fontenc}
\usepackage{hyperref}
\title{Wailaki Language Resources for K-12 Curriculum Development}
\author{Compiled by Claude Sonnet 5.5 (claude-sonnet-5-5) for Andr\'as Kornai}
\date{}
\begin{document}
\maketitle

\section{Introduction}
Wailaki (ISO 639-3 wlk) is an extinct Athabaskan language of northern California. The Berkeley Cal Language Archive page says no first-language speakers of any Eel River dialect remain in the 21st century. The search found no published dictionary, textbook or school programme. The material below is archival and academic.

\section{Primary Educational Resources}
No primer, textbook or school curriculum for Wailaki turned up in search. A heritage curriculum would have to be built from the archival and grammatical sources below.

\section{Linguistic and Reference Materials}
\begin{itemize}
\item Wikipedia, Wailaki language, with a short phonology and references: \url{https://en.wikipedia.org/wiki/Wailaki_language}
\item California Language Archive (UC Berkeley), Wailaki / Eel River Dene. The search excerpt shows it lists archival materials. I did not see the individual holdings. \url{https://cla.berkeley.edu/languages/eel-river-athabaskan.html}
\item Kayla Rae Begay, \emph{Wailaki Grammar}, PhD dissertation, UC Berkeley, 2017. This is cited in search results (the ASJP entry lists it as a source). I saw no direct URL for the dissertation.
\item Pliny Goddard, ``Wailaki Texts'' (1923, \emph{International Journal of American Linguistics}), and William Seaburg, ``A Wailaki (Athapaskan) Text with Comparative Notes'' (1977). Both are cited in the Wikipedia references. I saw no direct URLs.
\item ASJP word-list entry for Wailaki: \url{https://asjp.clld.org/languages/WAILAKI}
\end{itemize}

\section{Religious and Community Resources}
Nothing found in search. I did not see any Wailaki community organisation or scripture translation.

\section{Digital and Audio Resources}
No audio, video or software resources were found in search.

\section{Curriculum Development Context}
Wailaki is dormant, with no living speakers. A curriculum would be a reclamation effort based on Begay's grammar, Goddard's texts and archival recordings, ideally developed with Wailaki descendants and the Berkeley archive. Resources are very limited.
\end{document}